% Propietario conservado: /Users/ruben/Documents/New project/output/EXTREMA_MEDIA_RAZON_RECIPROCIDAD_ES_EN_20260918/ARTICULO/ES/source/sections/extension.tex
\section{Development of the formal kernel: initial conditions, geometry and refinement}
\label{x_reciprocidad:sec:extension}

The support of eighty-one positions does not coincide with the domain of initial conditions of the dynamics. An initial condition must also specify an orientation; and the additive and multiplicative sheets have independent cursors. This distinction allows emission to be reconstructed exhaustively without selecting a numerical region in advance.

\subsection{Initial states and the emission recurrence}

For enumeration we use indices \(I_9=\{0,\ldots,8\}\), related to the positive labels by \(i\mapsto i+1\). Let
\[
 \mathcal S=I_9^2\times\{N,E,S,O\},\qquad
 \Omega_{\rm seed}=\mathcal S_+\times\mathcal S_\times.
\]
Each initial condition \(\omega=(s_+,s_\times)\) determines two positions and two directions. The domain comprises all ordered pairs, so
\begin{equation}
 |\mathcal S|=324,\qquad |\Omega_{\rm seed}|=324^2=104\,976.
 \label{x_reciprocidad:eq:semillas}
\end{equation}
The census of residual-emitter signatures in Section~\ref{x_reciprocidad:app:firmas-emision} enumerates these conditions by identifiers and presents their fibers; the law transforming them is defined here. The rows of an emission table are fibers of this initial-condition domain.

The finite realization uses six windows of nine transitions. At each instant \(t\), the phase signature is \(M_{\rm ph}(1+((t-1)\bmod9))\). If it equals \(+1\), the positive residue \(\Sigma(i+1,j+1)=\rho_9(i+j+2)\) of the first cursor is read and that cursor is then shifted; if it equals \(-1\), the positive residue \(\Pi(i+1,j+1)=\rho_9((i+1)(j+1))\) of the second cursor is read and the second is then shifted. Zero increments the neutral counter. At the end of transitions twenty-seven and fifty-four, the cardinal orientations are reversed. Integer evaluations and their quotients are preserved in the register, but the emitter uses the residue readings specified here.

In the unit sector, the multiplicative reader uses the base-two discrete logarithm in \((\ZZ/9\ZZ)^\times\):
\[
 \ell_9(1)=0,\quad\ell_9(2)=1,\quad\ell_9(4)=2,\quad
 \ell_9(8)=3,\quad\ell_9(7)=4,\quad\ell_9(5)=5.
\]
For finite emission it is extended by zero to the nonunit sector. This extension is an observable, not an identification of the elements \(3,6,9\) with the multiplicative unit. The sheet and its membership in the radical sector remain in the trajectory register.

Let \(S_m\) be the sum of the three positive additive residues \(\Sigma\) of window \(m\), \(E_m\) the sum of the three discrete exponents of the multiplicative residues \(\Pi\) and \(Z_m=3\) the number of neutral events. Thus, \(S_m\) does not sum the integer evaluations \(S=i+j\), whose quotients remain in the register. With \([a]_{10}\in\{0,\ldots,9\}\), the decimal emission rule is
\begin{align}
 d_{m,1}&=[S_m]_{10},\notag\\
 d_{m,2}&=[S_m+E_m]_{10},\notag\\
 d_{m,3}&=[3S_m+5E_m+7Z_m]_{10},\notag\\
 U_m&=100d_{m,1}+10d_{m,2}+d_{m,3}.
 \label{x_reciprocidad:eq:emisor-seis}
\end{align}
The integer coefficients and calendar of this recurrence are an explicit part of the declared emission law. No values of \(\pi,\varphi,e\) or \(\alpha\) are received. The distinction between an emission rule and a constant recognized from its output prevents the genealogy from being hidden in a table of decimals.

If \(s_m=[S_m]_{10}\) and \(e_m=[E_m]_{10}\), the same rule is written
\begin{equation}
 U_m=100s_m+10[s_m+e_m]_{10}+[3s_m+5e_m+1]_{10}.
 \label{x_reciprocidad:eq:acoplamiento}
\end{equation}
The final one is the residue of \(7Z_m=21\) modulo ten. The oriented ternary reduction of the catalogue adopts the convention
\begin{equation}
 \Psi(U_1,\ldots,U_6)=([-U_1]_3,\ldots,[-U_6]_3).
 \label{x_reciprocidad:eq:psi-catalogo}
\end{equation}
The sign of this chart is explicit: changing it modifies the words and requires subsequent orientations to be transported. Balanced identification of each component of \(\FF_3\) with TRIT is performed after \eqref{x_reciprocidad:eq:psi-catalogo}.

\begin{proposition}[Injective factorization of emission]
The word \(U=(U_1,\ldots,U_6)\) determines the two signatures \(s=(s_1,\ldots,s_6)\) and \(e=(e_1,\ldots,e_6)\). The emission image has \(18\cdot26=468\) elements. The ternary projection has \(243\) distinct values in the ambient space of \(729\) words.
\end{proposition}
\begin{proof}
The hundreds digits of \(U_m\) recover \(s_m\). If \(b_m\) is its tens digit, then \(e_m=[b_m-s_m]_{10}\). The units digit checks the third congruence of \eqref{x_reciprocidad:eq:acoplamiento}. Coupling is therefore injective on the product of signatures.

Enumerating the \(324\) conditions for one cursor through the preceding recurrence produces eighteen additive signatures, each with eighteen preimages, and twenty-six multiplicative signatures, of which twenty-four have eight preimages, one has twenty-four, and one has one hundred and eight. Section~\ref{x_reciprocidad:app:firmas-emision} collects the signatures, their multiplicities, and the exact rule recovering their initial conditions. Injectivity of the coupling gives \(18\cdot26\) emissions and multiplicities
\[
 \underbrace{432}_{18\cdot24}\text{ fibers of }144,
 \qquad18\text{ fibers of }432,
 \qquad18\text{ fibers of }1944.
\]
The sum is \(432\cdot144+18\cdot432+18\cdot1944=104\,976\). Applying \eqref{x_reciprocidad:eq:psi-catalogo} to the \(468\) emissions produces exactly \(243\) words. This last number is the cardinality of the image, not of the space \(\FF_3^6\).
\end{proof}

This factorization explains the difference between finite exhaustiveness and unlimited depth. The domain \eqref{x_reciprocidad:eq:semillas} is complete for the declared support and emitter; it does not contain in advance every prefix of every prolongation. Depth belongs to the system of compatible extensions constructed on those conditions.

\subsection{Calendar closure and memory advance}

The calendar of one hundred and eight transitions consists of twelve nonadic windows. Its group structure is more precise than an informal product of two periods:
\begin{equation}
 0\longrightarrow C_{12}\xrightarrow{\iota}C_{108}
 \xrightarrow{p}C_9\longrightarrow0,
 \quad\iota([b])=[9b],\quad p([t])=[t]_9.
 \label{x_reciprocidad:eq:extension108}
\end{equation}
If \(t=9b+r\) with \(0\le r<9\), adding two phases produces the carry
\[
 (b,r)+(b',r')=
 \left(b+b'+\left\lfloor\frac{r+r'}9\right\rfloor,\ [r+r']_9\right),
\]
where the first coordinate is reduced modulo twelve. Composition preserves the term linking the two coordinates.

\begin{proposition}[Nonsplitting of the calendar]
The sequence \eqref{x_reciprocidad:eq:extension108} is exact and admits no section that is a group homomorphism.
\end{proposition}
\begin{proof}
The kernel of reduction modulo nine is the subgroup of multiples of nine, the image of \(\iota\). If the extension split, commutativity would give \(C_{108}\cong C_{12}\times C_9\). The first group has exponent \(108\), whereas the second has exponent \(\operatorname{mcm}(12,9)=36\), a contradiction.
\end{proof}

Nine transitions restore the nonadic phase and advance one window; one hundred and eight restore the finite calendar. Neither fact requires erasing the trajectory or restoring the prefix of a prolongation. The nonadic holonomy \(\Gamma_9\) denotes the transport of one turn on states with memory, so
\begin{equation}
 g(\Gamma_9x)=g(x),\qquad M(\Gamma_9x)=M(x)+1.
 \label{x_reciprocidad:eq:holonomia-memoria}
\end{equation}
Phase equality and memory increment are compatible and constitute two observations of the same transport.


% BEGIN OWNER /Users/ruben/Documents/New project/output/EXTREMA_MEDIA_RAZON_RECIPROCIDAD_ES_EN_20260918/ARTICULO/ES/source/sections/temporalidad_trit.tex
% Recuperación de 02_trit_geometrias y de la holonomía con memoria.
\subsubsection{Temporal parametrization of compatibility and tritic orientation}
\label{x_reciprocidad:trit-temporal-compatibilidad}

The temporal interpretation of TRIT is established on transitions of the
complete state. Let \(x\xrightarrow{\rm adm}y\) be the admissibility relation determined
by TPK operations on \(\widehat X\). A trajectory parametrized
by a discretely ordered set \(\mathcal T\) is a map
\[
 \gamma:\mathcal T\longrightarrow\widehat X,\qquad
 \gamma(t)\xrightarrow{\rm adm}\gamma(t^+),
\]
where \(t^+\) is the successor of \(t\). The compatibility relation precedes
choice of the parameter: the order allows its transitions to be expressed as
an evolution. When the term section is used, the projection
\(p:E\to\mathcal T\) and condition
\(p\circ\gamma=\operatorname{id}_{\mathcal T}\) are also declared. This specifies
the state space, temporal order and information transported.

The signature \(\tau\) selects a quadratic regime and admits a temporal
reading relative to that parametrization. Type \(+1\), with
\(J_{+1}^{2}=-I\), represents return and memory; type \(0\), with
\(J_0^{2}=0\), represents the transition threshold; type \(-1\), with
\(J_{-1}^{2}=I\), represents opening in two eigendirections. In the
model's temporal terminology, these functions correspond,
respectively, to past, present and future. The correspondence refers
to the operational position of a transition within a history:
preserved antecedent, update and admissible prolongation.

This interpretation acquires content through three distinct operations.
Restricting a trajectory recovers the already constructed prefix. The
update \(\mathcal U_t\) composes selection, transport and recording in
the observed transition. Prolongation considers the continuations that
preserve that prefix's boundary conditions. If
\(\mathcal P_n(x)\) denotes the set of admissible continuations of
length \(n\) from \(x\), removing the last step defines
\[
 r_n:\mathcal P_{n+1}(x)\longrightarrow\mathcal P_n(x).
\]
An unlimited continuation is a family \((\gamma_n)_{n\geq1}\)
satisfying \(r_n(\gamma_{n+1})=\gamma_n\). The inverse system thus preserves
the restrictions each depth imposes on the
preceding ones. Existence of a compatible family is established in the
corresponding prolongation domain; merely establishing the
sets \(\mathcal P_n(x)\) does not replace that proof.

Bidirectionality therefore combines reconstruction of antecedents
with compatibility of continuations. On the sequence of signatures, the
involution \(C_{\rm rev}\) reverses the order and changes the sign of each
component. The trajectory is additionally acted on by the position,
sheet and boundary transformations specified by transport. The inverted
signature is an observation of the conjugate trajectory, not a replacement
for its remaining coordinates. The threshold remains distinguished: the visible
value \(0\) retains its transition function even when its quotient
or historical register is nonzero.

The relation to exponential geometries is equally precise. In
each fixed regime, \(\exp(tJ_\tau)\) composes parameters and preserves the
algebraic norm. Elliptic closure belongs to that local representation.
The phase return of \(\Gamma_9\), by contrast, acts on the trajectory
with memory, according to \eqref{x_reciprocidad:eq:holonomia-memoria}. A trajectory can
recover its phase and preserve a new antecedent: the historical datum has
changed even though the phase observation is the same. The parabolic threshold
retains a linear shift; hyperbolic opening separates two
eigendirections of expansion and contraction. These three realizations
provide composition laws, while TPK determines their
selection and effective concatenation.

The aperiodic suspension of Section \ref{x_reciprocidad:sec:generacion} will make
this distinction explicit through a finite phase factor, an irrational rotation
and a memory degree. The link between TRIT and temporality is thus
expressed by operations on trajectories and their quadratic
representations. Its content goes beyond a classification of signs: it preserves the
regime, direction of transport and relation between antecedent,
transition and continuation.

% END OWNER


The monodromy representation of this return is developed in
Section~\ref{x_reciprocidad:mono-ampl:conexion}. There the holonomy of the
state is distinguished from its representation on phase and counter. The symbolic
suspension of Section~\ref{x_reciprocidad:mono-ampl:cristal} defines the
aperiodic time-crystal model, and Section~\ref{x_reciprocidad:mono-ampl:euler}
composes that temporality with the three Euler realizations.

\subsection{Observable compression and orthogonal conservation}

An isometric realization allows what a reduced observation loses to be expressed quantitatively. Let \(U:\mathcal H\to\mathcal H\) be an isometric realization of the transport of one turn, \(U^*U=I\), and \(J:\mathcal H_{\rm vis}\to\mathcal H\) an isometric inclusion. Define
\[
 T=J^*UJ,\qquad D=(I-JJ^*)UJ.
\]
Then
\begin{equation}
 UJ=JT+D,\qquad J^*D=0,\qquad I-T^*T=D^*D.
 \label{x_reciprocidad:eq:conservacion-ortogonal}
\end{equation}
The last equality follows by taking the adjoint product of the first and using orthogonality. Contraction of the visible component does not by itself constitute destruction of information in the total state.

\begin{proposition}[Conservation at a finite horizon]
If \(I-T^*T=D^*D\), then, for every \(n\ge1\),
\begin{equation}
 I=\sum_{r=0}^{n-1}T^{*r}D^*DT^r+T^{*n}T^n.
 \label{x_reciprocidad:eq:telescopia}
\end{equation}
\end{proposition}
\begin{proof}
Multiplying \(D^*D=I-T^*T\) by \(T^{*r}\) and \(T^r\) and summing produces a telescoping sum. The interior terms cancel and \(I-T^{*n}T^n\) remains.
\end{proof}

The last summand preserves the terminal state. Suppressing it before justifying a limit would alter the equality. This distinction is the operational content of preservation of unity in compression: the norm is distributed among the visible component, memory and terminal, without being replaced by a single residue.

\subsection{Local development: commutation, Lorentzian signature, and discrete curvature}
\label{x_reciprocidad:vii:sec:bloque-compresion}
\label{x_reciprocidad:iv:extension-curvatura-mixta}

The adjacent diagonal block with equal diagonals modulo nine is unique. Indeed,
\[
 k^2\equiv(k+1)^2\pmod9\quad\Longleftrightarrow\quad2k+1\equiv0\pmod9
\]
gives \(k=4\). Its visible matrix is
\[
 B_c=\begin{pmatrix}7&2\\2&7\end{pmatrix}=7I+2R,
 \quad R=\begin{pmatrix}0&1\\1&0\end{pmatrix},\quad R^2=I.
\]
With \(P_\pm=(I\pm R)/2\), we obtain \(B_c=9P_++5P_-\). The unity of stochastic normalization and preservation of an indefinite form are two realizations of this block:
\[
 \frac{B_c}{9}=P_++\frac59P_-,\qquad
 \frac{B_c}{\sqrt{45}}=\exp(\eta R),\quad\eta=\frac12\log\frac95.
\]
For \(Q=\operatorname{diag}(1,-1)\), the second matrix satisfies \(B_c^{\mathsf T}QB_c=45Q\). This yields a local realization in \(SO^+(1,1)\); proving it does not require introducing a four-dimensional spacetime beforehand.

Uniform averaging over the classes \(C_1=\{1,4,7\}\), \(C_2=\{2,5,8\}\), \(C_0=\{3,6,9\}\) gives
\[
 M_\Pi=\begin{pmatrix}4&5&6\\5&4&6\\6&6&9\end{pmatrix},\quad
 M_\Sigma=\begin{pmatrix}5&6&4\\6&4&5\\4&5&6\end{pmatrix}.
\]
The characteristic polynomial of \(M_\Pi\) is
\[
 (\lambda+1)((\lambda-9)^2-72),
\]
whence its quadratic form has signature \((2,1)\). The local spectral gap \(9-5=4\), coefficient \(2\), aggregate difference \(51-45=6\) and integrated defect \(9\cdot6=54\) describe distinct levels of the same arithmetic comparison. They are not interchangeable indices of the same operator.

Discrete curvature can be described directly before this spectral realization. In residues, let \(S(x,y)=x+y\), \(P(x,y)=xy\) and \(\Delta_x,\Delta_y\) be the forward differences. The link variables are potential differences, \(A_x^T=\Delta_xT\), \(A_y^T=\Delta_yT\), and we fix
\[
 F(T_x,T_y)=\Delta_x\Delta_yT_y-\Delta_y\Delta_xT_x.
\]
Since \(\Delta_xS=\Delta_yS=1\), \(\Delta_xP=y\), and \(\Delta_yP=x\),
\begin{equation}
 F(S,S)=F(P,P)=0,\qquad F(S,P)=+1,\quad F(P,S)=-1\pmod9.
\end{equation}
Sum and product, taken separately, induce flat connections; the ordered mixture changes the sign of plaquette curvature. On summing over a rectangle with sides \(a,b\), interior edges cancel and the oriented boundary sum equals \(\pm ab\) modulo nine. An exact relation thus appears between composition order and boundary orientation.

\subsection{Cubic completion and two notions of boundary}

The enlargement of the cubic support of side nine to that of side ten admits an exact positional stratification:
\begin{equation}
 10^3=9^3+3\cdot9^2+3\cdot9+1=729+243+27+1.
 \label{x_reciprocidad:eq:cubo}
\end{equation}
The first term corresponds to the three interior coordinates; the second, to one new coordinate; the third, to two; and the last, to the vertex of three new coordinates. The enlargement shell contains \(271\) positions. By contrast, the internal boundary of the discrete cube of side nine contains \(9^3-7^3=386\). The difference between these counts is a difference of domain, not a conflict of values.

Dividing \eqref{x_reciprocidad:eq:cubo} by \(1000\) gives a normalized partition of unity. Cube incidence distinguishes six faces, twelve edges and eight vertices. The subsequent link with Steiner hexads and octads requires the corresponding combinatorial maps; cube cardinalities do not replace them. This completion here provides the relation between bases \(729\) and \(1000\), while exceptional prolongation will be developed after numerical generation, the dodecaphasic register and the constructions of \(\alpha\).


% BEGIN OWNER /Users/ruben/Documents/New project/output/EXTREMA_MEDIA_RAZON_RECIPROCIDAD_ES_EN_20260918/ARTICULO/ES/source/sections/revision_emrh.tex
\subsubsection{Stratification of the completion and restoration of the apex}
\label{x_reciprocidad:sec:revision-emrh}

The completion admits a set-theoretic description making its
operation explicit. Let \(E_9\) be the set of nine residue classes, positively
represented by \(1,\ldots,9\), and let \(\ast\) be an additional element.
The symbol \(\ast\) is distinct from the zero class, already represented by \(9\).
Consequently, enlargement from \(E_9\) to
\(\widehat E=E_9\sqcup\{\ast\}\) increases cardinality from nine to ten
without modifying the internal modular identification. The cubic completion is
\(\widehat C=\widehat E^3\), and the initial cube is \(C=E_9^3\).

\begin{definition}[Strata of combinatorial codimension]
For \(0\leq r\leq3\), define
\[
 S_r=\{(x_1,x_2,x_3)\in\widehat E^3:
             \#\{j:x_j=\ast\}=r\}.
\]
The apex of the completion is \(a_\ast=(\ast,\ast,\ast)\), and the
punctured completion is \(\widehat C^{\circ}=\widehat C\setminus\{a_\ast\}\).
\end{definition}

\begin{proposition}[Exact decomposition of the completion]
The strata are disjoint and satisfy
\begin{align}
 \widehat C&=S_0\sqcup S_1\sqcup S_2\sqcup S_3,
 &|S_r|&=\binom3r9^{3-r},\label{x_reciprocidad:eq:revision-estratos}\\
 |\widehat C^{\circ}|&=729+243+27=999,
 &|\widehat C|&=999+1=1000.\label{x_reciprocidad:eq:revision-999}
\end{align}
The punctured enlargement contains \(270\) elements additional to the initial
cube; the complete enlargement contains \(271\).
\end{proposition}
\begin{proof}
An element of \(S_r\) is obtained by choosing the \(r\) coordinates
containing \(\ast\), in \(\binom3r\) ways, and assigning to the remaining ones
one of the nine classes of \(E_9\). Classification by the number of
additional coordinates is exhaustive and its classes are disjoint. The unique
element of \(S_3\) is \(a_\ast\). Removing it subtracts one unit from the total
cardinality; restoring it recovers the complete Cartesian product.
\end{proof}

The expression \emph{holographic extreme and mean ratio}, abbreviated EMRH in
the corpus, here denotes this articulation between the initial body,
enlargement strata and unit restitution. Its immediate content
is the exact partition of a finite product and its normalization. The number
\(999\) represents a determined object, the punctured completion;
\(1000\) represents its closure by restitution of the apex. The added
unit therefore has an explicit incidence residence.

\begin{figure}[htbp]
\centering
\begin{tikzpicture}[x=1cm,y=1cm,>=Stealth,font=\small]
\draw[->] (0,0)--(12.4,0);
\foreach \x/\a/\b in {0/729/{S_0},4.0/243/{S_1},7.5/27/{S_2},10.6/1/{S_3}}{
  \draw (\x,0.12)--(\x,-0.12);
  \node[above=5pt] at (\x,0) {\(\a\)};
  \node[below=5pt] at (\x,0) {\(\b\)};
}
\draw[decorate,decoration={brace,amplitude=5pt}] (0,1)--(7.7,1);
\node[above=8pt] at (3.85,1) {Punctured completion: \(999\)};
\draw[decorate,decoration={brace,mirror,amplitude=5pt}] (0,-1)--(10.8,-1);
\node[below=8pt] at (5.4,-1) {Full completion: \(1000=999+1\)};
\node[align=center] at (10.6,1.5) {Apex\\\((\ast,\ast,\ast)\)};
\end{tikzpicture}
\caption{Strata of the cubic completion. The arrangement represents
the set-theoretic decomposition, not a metric scale. The \(243\) elements
of \(S_1\) are three coordinate strata of \(81\) elements; the \(27\)
of \(S_2\) are three strata of \(9\). The apex completes the partition.}
\label{x_reciprocidad:fig:revision-emrh}
\end{figure}

\subsubsection{Preservation of unity and dimensional compatibility}

Normalization is not limited to dividing an isolated identity. In dimension
\(d\), the product \(\widehat E^d\) carries the uniform measure
\(\mu_d(\{x\})=10^{-d}\). The coordinate projection
\(p_d:\widehat E^{d+1}\to\widehat E^d\) removes the last coordinate.

\begin{proposition}[Projective compatibility of normalized measures]
For all \(d\geq1\),
\[
 (p_d)_*\mu_{d+1}=\mu_d,
 \qquad \mu_d(\widehat E^d)=1.
\]
The mass of the stratum with \(r\) additional coordinates is
\(\binom dr9^{d-r}/10^d\), and the sum of the masses of all strata
is one.
\end{proposition}
\begin{proof}
Each point of \(\widehat E^d\) has ten preimages. Their total mass is
\(10\,10^{-(d+1)}=10^{-d}\). Equality for any subset
is obtained by summing over its points. The stratified formula follows
from the same count as \eqref{x_reciprocidad:eq:revision-estratos}; the binomial expression
\((9+1)^d\) gives total normalization.
\end{proof}

In dimension three, removing the apex leaves mass \(999/1000\), and
restoring it adds \(1/1000\). Renormalizing before restitution
would define another measure: the distinction is operationally relevant.
In particular, conservation of probability mass, preservation of
transport registers and thermodynamic dissipation are assertions about
different objects. The proposition proves the first and provides the
compatible support for studying the second; the third requires a
specified dynamic and energetic realization.

The internal geometric boundary also remains distinct.
On the Cartesian representative \(\{1,\ldots,9\}^3\), removing
\(\{2,\ldots,8\}^3\) leaves \(386\) points. This boundary belongs to the
initial cube, whereas \(\widehat C\setminus C\) belongs to its
enlargement. Cubic completion, base-one-thousand positional expansion and
the exceptional incidence reading share support data;
their respective maps determine what information from those data each
representation uses.

% END OWNER



% BEGIN OWNER /Users/ruben/Documents/New project/output/EXTREMA_MEDIA_RAZON_RECIPROCIDAD_ES_EN_20260918/ARTICULO/ES/source/sections/geometria_modular.tex
\subsection{Cubic incidence and Lorentzian realization}

Cubic completion provides a relation between interior and boundary.
Its geometry can be developed through incidence of the six faces
with the eight vertices. Label the faces by \((i,\varepsilon)\),
\(i\in\{1,2,3\}\), \(\varepsilon\in\{+1,-1\}\), and the vertices by
\(\nu\in\{+1,-1\}^3\). The incidence matrix is
\[
 M_{(i,\varepsilon),\nu}=\mathbf1_{\{\nu_i=\varepsilon\}}.
\]
Let \(u=(1,\ldots,1)^{\mathsf T}\in\RR^6\), \(R_F\) be face
opposition and \(P_u=uu^{\mathsf T}/6\), \(P_-=(I-R_F)/2\).

\begin{proposition}[Four-dimensional incidence quotient]
The incidence Gram matrix satisfies
\[
 MM^{\mathsf T}=12P_u+4P_-.
\]
Its eigenvalues \(12,4,0\) have multiplicities \(1,3,2\).
Consequently,
\[
 Q_\square=\RR^6/\ker(MM^{\mathsf T})
 \simeq\RR u\oplus E_-,\qquad E_-=\operatorname{im}P_-.
\]
\end{proposition}
\begin{proof}
A face contains four vertices; two opposite faces have empty
intersection, and two distinct nonopposite faces share two vertices.
The matrix \(12P_u+4P_-\) has precisely those entries.
The uniform space has dimension one; differences between opposite
faces span a three-dimensional space orthogonal to the uniform one.
Its complement has dimension two and belongs to the kernel.
\end{proof}

The decomposition \(1+3\) comes from incidence. The temporal choice
is declared through the bilinear form
\[
 B_\square([x],[y])=\langle x,(P_--P_u)y\rangle.
\]
It is well-defined on the quotient, because both projectors annihilate the
kernel. In the basis
\[
 t=[u]/\sqrt6,\qquad
 d_i=[e_{i,+}-e_{i,-}]/\sqrt2,
\]
its matrix is \(\operatorname{diag}(-1,1,1,1)\). The domain,
decomposition and sign convention producing
the Minkowski realization are made explicit. The positive Gram matrix determines the quotient;
the temporal assignment determines the indefinite form on it.

The operator
\[
 W_\square=\sqrt6P_u+\sqrt2P_-,\qquad
 W_\square e_{i,\varepsilon}=t+\varepsilon d_i
\]
sends the six face labels to null directions, since
\(B_\square(t+\varepsilon d_i,t+\varepsilon d_i)=-1+1=0\).
Inversion of a direction and face opposition thus have
a specific geometric realization.

\subsubsection{Soldering, inversion and subdivision}

The cellular realization uses an invertible transport \(U_m(a)\),
preserving the Lorentzian form, between the fibers at the endpoints of
an oriented dual edge \(a\), a face
label \(\lambda_m(a)\) and an orientation \(\sigma_m(a)\).
With \(\ell_m=\ell_0/9^m\), the soldering map is
\[
 \beta_m(a)=\ell_m\sigma_m(a)
 W_{\square,s(a)}e_{\lambda_m(a)}.
\]
Compatible reversal satisfies
\[
 \beta_m(\bar a)=-U_m(a)\beta_m(a).
\]
Identify the origin fibers at two consecutive levels through
the declared refinement isometry. If the nine subsegments
\(a_1,\ldots,a_9\), of level \(m+1\), preserve the label
and orientation when transported to that common fiber, then
\begin{equation}
 \beta_m(a)=\sum_{j=1}^9
 U_{m+1}(a_1\cdots a_{j-1})^{-1}\beta_{m+1}(a_j).
 \label{x_reciprocidad:eq:soldadura-refinamiento}
\end{equation}
Each transported term equals \(\ell_m/9\) times the same
null direction, proving the identity. The compatibility
hypothesis determines which subdivisions realize this refinement.
The equation combines direction, scale and transport, instead of limiting
the geometric correspondence to a count of dimensions.

With orientation and Lorentzian form also fixed, Hodge
duality on \(\Lambda^2Q_\square^*\) satisfies \(\star^2=-I\).
In dimension four, the sign is \((-1)^{2(4-2)+1}=-1\).
This yields a real complex structure on two-forms and,
after complexification, the eigenspaces of values \(+i\) and
\(-i\). The electromagnetic realization will use this domain
after constructing the electronic section and declaring its coupling.

\subsection{Compatibility between modular and positional reductions}

The block decimal base enters through a Euclidean division
simultaneously preserving remainder and quotient:
\[
 n=1000q+r,\qquad 0\le r<1000.
\]
Since \(1000\equiv1\pmod9\) and \(1000\equiv1\pmod3\),
\[
 n\equiv q+r\pmod9,\qquad n\equiv q+r\pmod3.
\]
Therefore, reduction of a block requires the quotient's
contribution to recover the integer's class. The numbers \(0\) and
\(1000\) have the same remainder modulo \(1000\), but different classes
modulo \(9\) and modulo \(3\). This is an operator reason to preserve
memory during a change of resolution.

Ternary refinement and determination of a decimal prefix
can be coordinated through a single integer residue. If \(N/3^L\)
is a cylinder endpoint and \(D/1000^r\) a decimal endpoint, define
\[
 A=1000^rN-D3^L.
\]
On adding six ternary symbols, \(N'=729N+\nu\),
\(L'=L+6\), with \(0\le\nu<729\). If the decimal prefix remains
fixed, we obtain
\[
 A'=729A+\nu1000^r.
\]
If a decimal block \(\delta\) is also incorporated, then
\(D'=1000D+\delta\), \(r'=r+1\), and
\begin{equation}
 A'=729000A+\nu1000^{r+1}-729\delta3^L.
 \label{x_reciprocidad:eq:residuo-dos-bases}
\end{equation}
Both identities are proved by substituting the definitions. Endpoint
comparison is thus reduced to integer operations retaining
the contribution of each prolongation. This compatibility
allows a region to be refined before its real coordinate is evaluated.

% END OWNER


% BEGIN OWNER /Users/ruben/Documents/New project/output/EXTREMA_MEDIA_RAZON_RECIPROCIDAD_ES_EN_20260918/ARTICULO/ES/source/sections/aritmetica_historias.tex
% Incorporación al final de extension.tex; no contiene una sección principal.
% Procedencia: RESULTADO_RECUPERADO; explicitación sectorial de las hipótesis.
% Fuentes leídas: monografía Doble proyección holográfica (cantera editorial),
% manuscrito/local/01_tesis_recuperada.tex y
% manuscrito/generated/algebra_holonomica.tex, especialmente §§1--7 y 21.
% Definición rectora de número: manuscrito/flattened/main_autosuficiente.tex,
% secciones que comienzan en líneas 22137 y 23000 de la cantera de 775 páginas.
% No se adopta su presentación global como prueba de todas las flechas TPK.

\subsection{The number as a compatible section and the value as a reading}
\label{x_reciprocidad:hist:seccion}

Distinguishing state from observation specifies what object is prolonged
as resolution increases. At depth \(n\), let \(X_n\) be the set of
admissible states produced by APP--TRIT--TPK. Their data include position,
sum and product sheets, residues and quotients, regime, orientation,
carry, phase, route, boundary, memory, and survival. These are not
independent coordinates: earlier arithmetic evaluations and transports
impose their relations. A prolongation preserves these relations in
each antecedent, even while adding new steps and memory data.
This is the meaning of the arithmetic of histories developed in
(Sections~\ref{x_reciprocidad:sec:extension} and~\ref{x_reciprocidad:mono-ampl:conexion}).

Let \(r_{m,n}:X_m\to X_n\), for \(m\geq n\), be the depth
restrictions, with \(r_{n,n}=\operatorname{id}\) and
\(r_{\ell,n}=r_{m,n}\circ r_{\ell,m}\). Each restriction preserves the
prefix and the information already determined in it.

\begin{definition}[Compatible numerical section]
An HMT number is a compatible history of the state system:
\begin{equation}
 x=(x_n)_{n\geq0}\in X_\infty:=\varprojlim_n X_n,
 \qquad r_{m,n}(x_m)=x_n.
 \label{x_reciprocidad:hist:limite-inverso}
\end{equation}
A reader is an additional operation on a specified domain of these
histories. A digit is a finite observation; a scalar value is an
evaluation of the reader. The pair formed by the section and its reader constitutes
a read presentation, not a redefinition of the section's identity.
\end{definition}

In the Archimedean domain \(X_\infty^{\mathrm{Arch}}\), the reader associates
a closed cylinder with each prefix, preserves their nesting and makes
their diameter tend to zero. It thus determines
\begin{equation}
 \operatorname{ev}_{\mathbb R}:X_\infty^{\mathrm{Arch}}\longrightarrow
 \mathbb R,\qquad
 \{\operatorname{ev}_{\mathbb R}(x)\}=\bigcap_n I_n(x).
 \label{x_reciprocidad:hist:evaluacion}
\end{equation}
The specific construction of these cylinders and determination of their boundaries
are developed in Section~\ref{x_reciprocidad:sec:generacion}. Definition
\eqref{x_reciprocidad:hist:limite-inverso} does not by itself assert that the image of
\eqref{x_reciprocidad:hist:evaluacion} is the whole line: a coverage assertion
requires its own theorem. Nor does equality of values establish,
without an injectivity result, equality of their histories. This separation
allows precision to be discussed as depth while preserving the difference
between the arithmetic antecedent and its scalar realization.

\subsection{Composition of histories and memory multiplier}
\label{x_reciprocidad:hist:algebra}

Fix a sector of states at depth \(n\) and a category
\(\mathcal C\) whose arrows are its admissible histories. Composition
is written \(vu=v\circ u\): first \(u\) is traversed, then \(v\).
Elementary steps come from TPK selection, transport and
update. A local tritic fiber admits the expression
\(\mathbb R[J_x]/(J_x^2+\tau(x)1_x)\), where
\(\tau(x)\in\{+1,0,-1\}\). Regime and orientation belong to the
transported state; a transition between regimes is not equivalent to
identifying their quadratic algebras.

The following formulation makes an algebraic sector of that composition precise.
For each object \(x\), a unital associative real algebra \(A_x\) is given,
which may be noncommutative. For each \(u:x\to y\), a
unital homomorphism \(\rho_u:A_x\to A_y\) is given. We assume a strict action:
\begin{equation}
 \rho_{\operatorname{id}_x}=\operatorname{id}_{A_x},\qquad
 \rho_v\rho_u=\rho_{vu}.
 \label{x_reciprocidad:hist:accion-estricta}
\end{equation}
For each composable pair \(u:x\to y\), \(v:y\to z\), fix an
invertible central multiplier
\(\Omega(v,u)\in Z(A_z)^\times\). We require the normalization
\(\Omega(\operatorname{id}_y,u)=\Omega(u,\operatorname{id}_x)=1_y\)
and, for each composable triple, the identity
\begin{equation}
 \rho_w\bigl(\Omega(v,u)\bigr)\Omega(w,vu)
   =\Omega(w,v)\Omega(wv,u).
 \label{x_reciprocidad:hist:cociclo}
\end{equation}
These are explicit hypotheses on transports and their coefficients, not a
consequence of calling the dynamics holonomic. When applying this presentation
to a TPK sector, its \(\rho_u\) and \(\Omega(v,u)\) must be identified.
Transitions expressed through bimodules require that structure
to be preserved and are not replaced by \eqref{x_reciprocidad:hist:accion-estricta}.

The space of finite sums
\[
 \mathcal H(\mathcal C,A,\rho,\Omega)
   =\bigoplus_{u\in\operatorname{Mor}(\mathcal C)} A_{t(u)}T_u
\]
receives the bilinear product
\begin{equation}
 (aT_v)(bT_u)=a\rho_v(b)\Omega(v,u)T_{vu},
 \label{x_reciprocidad:hist:producto}
\end{equation}
if the arrows are composable, and zero otherwise. Here \(t(u)\)
denotes the terminal state. The sum combines histories with coefficients; the
product concatenates histories while preserving the register multiplier.

\begin{proposition}[Sectorial associativity with a central multiplier]
Under \eqref{x_reciprocidad:hist:accion-estricta} and \eqref{x_reciprocidad:hist:cociclo}, the product
\eqref{x_reciprocidad:hist:producto} is associative. If the object set is finite,
its unit is \(\sum_x1_xT_{\operatorname{id}_x}\).
\end{proposition}
\begin{proof}
For \(u:x\to y\), \(v:y\to z\), \(w:z\to t\), let
\(a\in A_y\), \(b\in A_z\), \(c\in A_t\). The two coefficients
of \(T_{wvu}\) are
\begin{align*}
 ((cT_w)(bT_v))(aT_u):\quad&
 c\rho_w(b)\Omega(w,v)\rho_{wv}(a)\Omega(wv,u),\\
 (cT_w)((bT_v)(aT_u)):\quad&
 c\rho_w(b)\rho_w\rho_v(a)\rho_w(\Omega(v,u))\Omega(w,vu).
\end{align*}
The strict action equates \(\rho_w\rho_v(a)\) with \(\rho_{wv}(a)\).
Centrality of \(\Omega(w,v)\) allows it to be moved to the right of
that factor; \eqref{x_reciprocidad:hist:cociclo} then equates the remaining products.
The coefficients \(a,b,c\) are not permuted. If the triple is not composable,
both parenthesizations are zero because of their initial and terminal states.
Bilinearity completes the proof. The local identities and
normalization of \(\Omega\) verify the indicated unit.
\end{proof}

Calendar carry provides an effective example. In the category
with one object and arrows \(C_9\), take coefficients
\(A=\mathbb R[z,z^{-1}]\), identity action and
\[
 \Omega(b,a)=z^{c_0(a,b)},\qquad
 c_0(a,b)=\left\lfloor\frac{a+b}{9}\right\rfloor,
 \quad 0\leq a,b<9.
\]
The two carry sums of a triple are
\(\lfloor(a+b+c)/9\rfloor\); hence
\eqref{x_reciprocidad:hist:cociclo} holds. The formal variable \(z\) records the turn without
assigning it a numerical value. Subsequently imposing \(z^{12}=1\) recovers the
dodecaphasic register of the calendar \eqref{x_reciprocidad:eq:extension108}; retaining
\(z\) without that relation distinguishes successive turns. The example realizes
the memory of this calendar, without identifying it with the entirety of TPK
memory.

\subsection{Double variance of refinement and preservation of unity}
\label{x_reciprocidad:hist:doble-varianza}

Restricting states prolongs observables in the opposite direction.
Consider the function algebras, with pointwise product,
\(D_n=\operatorname{Fun}(X_n,\mathbb R)\). Precomposition defines
\begin{equation}
 r_{m,n}^*:D_n\longrightarrow D_m,
 \qquad r_{m,n}^*f=f\circ r_{m,n}.
 \label{x_reciprocidad:hist:precomposicion}
\end{equation}
These are unital morphisms; they are also injective when the corresponding
restrictions are surjective. The direct limit
\(D_{\mathrm{cyl}}=\varinjlim_n(D_n,r_{n+1,n}^*)\) brings together
observations of finite depth. This algebra is not identified here
with the entire route algebra: prolonging transport operators also requires
their compatibilities with refinement.

\begin{lemma}[Unit and naturality of evaluation]
If \(e_x\) is the indicator of \(x\in X_n\), then
\begin{equation}
 r_{m,n}^*e_x=\sum_{y:\,r_{m,n}(y)=x}e_y,
 \qquad r_{m,n}^*1_n=1_m,
 \label{x_reciprocidad:hist:unidad}
\end{equation}
and, for every \(f\in D_n\), \(y\in X_m\),
\[
 \langle r_{m,n}^*f,y\rangle=\langle f,r_{m,n}y\rangle.
\]
Each section \(x\in X_\infty\) determines a well-defined evaluation
\(D_{\mathrm{cyl}}\to\mathbb R\), given by \([f]\mapsto f(x_n)\)
for \(f\in D_n\).
\end{lemma}
\begin{proof}
The fibers of \(r_{m,n}\) partition \(X_m\). Their
indicators are orthogonal and sum to \(1_m\), proving
\eqref{x_reciprocidad:hist:unidad}. Naturality is the definition of precomposition.
Finally, \((r_{m,n}^*f)(x_m)=f(x_n)\), so two compatible
representatives give the same evaluation in the direct limit.
\end{proof}

The global unit is the constant function one over all possibilities;
the individual section selects a compatible history. Preserving the
former does not mean immobilizing the latter. Refinement may increase
the number of extensions and the depth of each history while maintaining
unity and all evaluations of its antecedents.

\subsection{Phase return, Weyl representation and two readings}
\label{x_reciprocidad:hist:puente-lecturas}

The law \eqref{x_reciprocidad:eq:holonomia-memoria} acquires an operator realization
in the bilateral suspension to be constructed in
Section~\ref{x_reciprocidad:sec:generacion}. There \(T\) is the one-step advance,
\(V_9\) observes the nonadic phase, \(W_\delta\) the rotation phase and
\(D_M\) the memory degree. Writing
\(\widehat\Gamma_9=T^9\) for the turn in this representation,
the Weyl relations \eqref{x_reciprocidad:eq:weyl-relojes} imply, on finitely supported
vectors,
\begin{equation}
 \begin{aligned}
 V_9\widehat\Gamma_9&=\widehat\Gamma_9V_9,\qquad
 W_\delta\widehat\Gamma_9=e^{2\pi i9\delta}
     \widehat\Gamma_9W_\delta,\\
 [D_M,\widehat\Gamma_9]&=\widehat\Gamma_9.
 \end{aligned}
 \label{x_reciprocidad:hist:tres-retornos}
\end{equation}
The first equality follows from \(\zeta_9^9=1\); the second, from nine
applications of the covariance of \(W_\delta\); the third expresses
the unit memory increment. Finite phase returns, while
irrational rotation and memory distinguish the next state. The
characters and parameter \(\delta\) are subsequently realized from the
numerical construction specified there: they do not retrospectively select
the seeds. This representation describes observables and translations of
the suspension; it does not replace all coordinates or all operators
of TPK.

The double reading prolongs the same distinction between history and observable.
In its Archimedean domain, \eqref{x_reciprocidad:hist:evaluacion} is applied; in the
calibrated incidence domain, the maps \(Q_n\) of
Section~\ref{x_reciprocidad:exc:doble-lectura} preserve the words and boundary
frames. The equality proved there
\(r_{n,m}Q_n=Q_mp_{n,m}\) determines their prolongation \(Q_\infty\).
On the common domain of both readings one may jointly consider
\(\operatorname{ev}_{\mathbb R}(x)\) and \(Q_\infty(x)\): the
first retains the value; the second, the transported incidence preserved by its
codomain. The section precedes both. This articulation allows the passage
from arithmetic of histories to constants and the exceptional
reading without turning their realizations into new initial data.

% END OWNER

